\documentclass[11pt,a4paper]{article}
\usepackage[T1]{fontenc}
\usepackage[utf8]{inputenc}
\usepackage[english]{babel}
\usepackage{amsmath,amssymb}
\usepackage[margin=2cm]{geometry}
\setlength{\parindent}{0pt}
\setlength{\emergencystretch}{3em}
\allowdisplaybreaks

\title{About \texttt{quantified\_inductions.tex}}
\author{}
\date{}
\begin{document}
\maketitle
This note explains what the test data in \texttt{quantified\_inductions.tex} is used for. That file is input for the engine in \texttt{main.c}; it is read by the test \texttt{test\_simple\_inductions.c} in this folder.

\begin{enumerate}
\item \textbf{Purpose.} The file holds 14 true statements in which a formula with a sum or a product is stated ``for every $n$ in a set of numbers'', with the quantifier written out. The engine must prove every one of them.

\item \textbf{What it exercises.} A statement with a quantifier goes to the rules of chapter 18 of \texttt{main.c}, which reason about polynomials. A sum up to $n$ is not a polynomial, so those rules do not apply: the engine has to recognise that a universal statement over the natural numbers, about an equality, can be proved by induction, and call the method of chapter 13 from inside the other one. An earlier version answered that only rational expressions are allowed.

\item \textbf{Variety.} The starting point is $0$, $1$ or $2$ according to the set; the letter is not always $n$; two statements are conjunctions, one of them of an induction and of a statement about rational numbers; one asks for a value of $n$ instead of every $n$.

\item \textbf{Expected result.} The line \texttt{PASS: 14 of 14 statements proved, each with a written proof}.

\item \textbf{How to run it.} From the top folder of the repository:
\begin{verbatim}
cc -std=c11 -O2 -Wall -Wextra -o /tmp/test_simple_inductions \
    tests/sums_products_inequalities/test_simple_inductions.c \
    && /tmp/test_simple_inductions
\end{verbatim}
\end{enumerate}
\end{document}
